\section{The zones with the most group flights}\label{sec:zones}

Do the three populations differ in the same way where groups are most frequent?

A zone is a cell of \RegionCell~km $\times$ \RegionCell~km that holds the middle of a piece, whichever
the takeoff of the flight; the pilots of a zone need not take off from the same place. The
\RegionZones{} cells with the most flights that have group time in them are kept, and each is named
by the takeoff most flown by its group flights (the cells are on a fixed grid; Figure~\ref{fig:maps}). In each zone the pieces of $\Sl$, $\G$ and $\A$ are those of
Section~\ref{sec:populations}, the measure is the same, and a population is drawn when it has at least
\RegionMinFlights{} flights in the zone. Bands are 5--95\% of bootstrap draws that resample the dependency
blocks of the grouping.

\begin{figure}
\centering
\includegraphics[width=\textwidth]{regions.pdf}
\caption{$\Hloc(\tau)$ of $\Sl$, $\G$ and $\A$ in the \RegionZones{} zones, with the 5--95\% band and,
inset, $\MSD(\tau)$ in $\mathrm{m}^2$ (log--log). The number of flights of each population is in the
legend and in Table~\ref{tab:zones}.}
\label{fig:regions}
\end{figure}

In \RegionOrdered{} of the \RegionZones{} zones $\Hloc(1000~\mathrm{s})$ is highest in group time, then
in the same flights alone, then in solo flights; in the others the same flights alone are as straight as
group time or more (Figure~\ref{fig:regions}, Table~\ref{tab:zones}). Each zone holds few flights of
group time.

\begin{table}
\centering\footnotesize
\begin{tabular}{@{}l r r r r c c c@{}}
\toprule
zone & group flights & flights $\Sl$ & flights $\G$ & flights $\A$ & $\Hloc(1000~\mathrm{s})$ $\Sl$ & $\G$ & $\A$ \\
\midrule
St Marc d'Ouilly (49.0$^\circ$N, -0.4$^\circ$E) & 247 & 628 & 247 & 103 & 0.939 & 0.944 & 0.962 \\
les richards (45.4$^\circ$N, 6.1$^\circ$E) & 179 & 414 & 179 & 210 & 0.917 & 0.942 & 0.919 \\
Douelle (44.5$^\circ$N, 1.5$^\circ$E) & 157 & 392 & 157 & 56 & 0.938 & 0.943 & 0.953 \\
La Mounine (44.5$^\circ$N, 2.2$^\circ$E) & 154 & 248 & 154 & 86 & 0.946 & 0.946 & 0.947 \\
Soult (46) (44.9$^\circ$N, 1.5$^\circ$E) & 142 & 323 & 142 & 78 & 0.946 & 0.952 & 0.951 \\
Jeufosse (49.0$^\circ$N, 1.5$^\circ$E) & 135 & 214 & 135 & 52 & 0.947 & 0.974 & 0.963 \\
\bottomrule
\end{tabular}

\caption{The zones, by the number of flights with group time in them. Flights of each population and
$\Hloc(1000~\mathrm{s})$; the bands are in Figure~\ref{fig:regions}.}
\label{tab:zones}
\end{table}

\paragraph{What we keep.} Zone by zone the three populations are close, and their order changes from one
zone to the other.
